\documentclass[11pt,a4paper]{article}
\usepackage[T1]{fontenc}
\usepackage{lmodern}
\usepackage[margin=24mm,headheight=15pt]{geometry}
\usepackage{microtype}
\usepackage{amsmath,amssymb,booktabs,tabularx}
\usepackage{xcolor,fancyhdr}
\usepackage[colorlinks=true,linkcolor=teal,urlcolor=teal]{hyperref}

\definecolor{ink}{HTML}{183544}
\definecolor{teal}{HTML}{08717C}
\hypersetup{
  pdftitle={AniMoViz Similarity Methodology},
  pdfauthor={AniMoViz project},
  pdfsubject={Catalogue sampling, Qwen3 embeddings, and cosine similarity}
}
\pagestyle{fancy}
\fancyhf{}
\fancyhead[L]{\small\textcolor{ink}{AniMoViz}}
\fancyhead[R]{\small Similarity methodology}
\fancyfoot[R]{\small\thepage}
\renewcommand{\headrulewidth}{0.3pt}
\setlength{\parindent}{0pt}
\setlength{\parskip}{6pt}
\renewcommand{\arraystretch}{1.21}
\newcommand{\file}[1]{\texttt{\detokenize{#1}}}
\newcommand{\flag}[1]{\texttt{-\kern0pt-\detokenize{#1}}}
\newcommand{\heading}[1]{\par\bigskip{\large\bfseries\textcolor{ink}{#1}}\par\nobreak\smallskip}

\begin{document}

{\Huge\bfseries\textcolor{ink}{AniMoViz similarity methodology}}\par
\medskip
AniMoViz compares the titles, genres, tags, and synopses of 1,000 anime
and 1,000 movies. The resulting one million pair scores support the
Dust \& Magnet view and the annual matches in History.

\heading{How titles enter the catalogues}

The shared catalogue collects 1,000 unique non-adult anime in AniList
popularity order across formats and release years. It scans the TMDB
popular-movie list until it has 1,000 unique non-adult movies. The sample
reflects popularity at the time of collection.

When both sources list the same version of a film, the AniList record
stays and the next eligible movie fills the TMDB slot. The refresh uses
reviewed ID pairs and checks titles and aliases, format, release dates,
runtimes, and synopsis evidence. Alternate cuts and adaptations can keep
separate entries. The staged \file{duplicate_report.json} records
exclusions and ambiguous pairs kept for review.

\heading{What goes into each embedding}

Each title becomes one labeled document using these fields.

\begin{tabularx}{\linewidth}{@{}p{0.17\linewidth}X@{}}
\toprule
Field & Preparation \\
\midrule
Title & Use the prepared title. If it is empty or visibly damaged, try the English title, then romaji. \\
Genres & Clean and deduplicate the genre labels within the field. \\
Tags & Clean and deduplicate AniList tags or TMDB keywords. AniList tags marked as spoilers are excluded during collection. \\
Synopsis & Clean the plot description and preserve its wording. \\
\bottomrule
\end{tabularx}

The scorer decodes HTML entities, removes markup, normalizes Unicode,
and collapses whitespace. It omits empty fields and retains records
with usable text. Ratings, popularity, IDs, release dates, alternate titles,
and credits remain outside the embedding input.

\heading{Model and inference}

Both catalogues use the same local \href{https://huggingface.co/Qwen/Qwen3-Embedding-4B}{Qwen3-Embedding-4B}
model with an empty prompt. The model pools the last non-padding
token into a 2,560-dimensional vector. Inference runs in float16.
The scorer saves normalized float32 vectors and records its input hashes,
model revision, runtime versions, and settings. The pinned revision is
\file{5cf2132abc99cad020ac570b19d031efec650f2b}.

The default budget is 2,048 tokens including special tokens. The scorer
checks the budget before inference. A standalone run can raise the budget
or explicitly permit right-side truncation, which the record map marks.

\newpage

{\LARGE\bfseries\textcolor{ink}{From vectors to the three views}}\par

\heading{One matrix for both directions}

Let \(a_i\) be the anime vector and \(m_j\) the movie vector. Cosine
similarity is the dot product of their unit vectors.
\[
  s_{ij}
  = \frac{a_i^{\mathsf T}m_j}{\lVert a_i\rVert_2\,\lVert m_j\rVert_2}
  = \widehat a_i^{\mathsf T}\widehat m_j,
  \qquad -1 \leq s_{ij} \leq 1 .
\]
Stacking the unit vectors as rows gives \(S=\widehat A\widehat M^{\mathsf T}\)
with 1,000 anime rows and 1,000 movie columns. The scorer clamps floating
point excursions to the cosine range and preserves negative scores.
Both directions use the same pair score. The research array follows
the CSV orders. The web export stores the same values movie-first in
\file{movie_anime_scores.bin} as little-endian float32. Its companion
\file{movie_recommendations.json} gives the axis order, matrix hash, and the
20 highest and 20 lowest anime matches per movie.

\heading{How the views use the data}

\textbf{Dust \& Magnet.} A title magnet uses its pairwise cosines.
A cross-catalogue genre magnet averages signed cosines over sampled
titles in the genre. Studio and person groups average across their
credited works for similarity magnets. When a popularity, rating, or
count magnet is active, the view converts each magnet's values to
midrank percentiles using dust with data for all active magnets.
The field takes an equal share of strong matches for each magnet,
then fills remaining slots by mean cosine or attraction percentile.

\textbf{History.} The annual AniList archive samples up to ten Japanese,
non-adult, released anime per first-release year from 1963 onward.
The Anime Timeline takes the title with the most list users each year
and finds up to five movie matches by cosine. Eligible movies have a
release year and genre metadata and exclude Animation and the
keyword \textquotedblleft cartoon characters.\textquotedblright{}
Anime outside the shared matrix are scored with the same model and
movie corpus. The Movie Timeline starts with the highest-popularity
eligible TMDB movie for each year and uses its ranked anime matches
from the shared catalogue.

\textbf{Full Catalogue.} Search, filters, and side-by-side comparisons use
metadata for the same two samples. Its ratings and popularity come from
the source snapshots.

\heading{Interpretation and reproducibility}

Higher cosine values indicate closer text representations. A cosine of
0.8 has no percentage or probability interpretation. Shared genre labels,
franchise entries, incomplete synopses, and differences in source tags can
influence rankings. The popularity sample favors widely listed works.

Use \file{data/refresh_data.py} with \flag{update-dataset} to stage a new
catalogue. Run it with \flag{recompute-scores} to score and export it.
The completed run saves source-to-row maps, model settings, and checksums.
Library and hardware changes can introduce small numerical differences.

\end{document}
